\chapter{Literature Review}

Since my thesis project is predominantly engineering, the literature
review mostly focuses on existing practical solutions for i18n in
software products. Although scientific research on this topic has
been conducted, it is rather sparse. Therefore, to avoid being
out-of-touch with the state of the art, my main priority is to
analyze and compare real-world technologies.

\section{Related Research}

\textit{
  Более подробно рассказать про \cite{barbour1997, he2002,
  couto2026}, их находки и выводы, что я могу применить в своём
  проекте, с чем расхожусь, что не беру в область своей работы.
}

\section{Existing Projects}

\textit{
  Сделать сравнительный анализ и описание существующих проектов
  и чем я буду выделяться на фоне них, от чего вдохновляюсь, что
  переиспользую, с кем конкурирую и какие у моего решения преимущества.
}

ICU MessageFormat \cite{icumf}, MessageFormat 2 \cite{icumf2},
CLDR \cite{cldr}, ICU \cite{icu}, ICU4X \cite{icu4x}

\cite{typesafe-i18n} \cite{rust-i18n} \cite{typed-i18n} \cite{rosetta}

\textbf{gettext} \cite{gettext} is a traditional i18n framework with
very widespread adoption. Virtually every programming language has
libraries for gettext support; however, only two claim to be type-safe:
i18n++ for C++ \cite{i18n-cpp} (unmaintained, low adoption,
type-safety through macros) and i18n-embed for Rust \cite{i18n-embed}
(type-safety through macros). With gettext, programmers write
user-facing text messages in the main locale (most often, English)
directly in their code, and translation files are used to map these
strings to other locales' versions. The framework uses custom syntax
with support for dynamic placeholders and pluralization. However,
more advanced features such as date/time/number formatting, message
variants, or markup placeholders are not supported.

\textbf{i18next} \cite{i18next}...

\textbf{Paraglide JS} \cite{paraglide}...

\textbf{Lingui} \cite{lingui}...

\textbf{doti18n} \cite{doti18n}...

\textbf{Project Fluent} \cite{fluent}...

Fluent-Typed \cite{fluent} i18n-embed \cite{i18n-embed}

\textbf{Mulan}

\begin{longtable}{c|c|c|c|c}
  \caption{Rating Existing i18n Solutions} \label{tab:pfams} \\
  \hline
  \textbf{Name} &
  \textbf{Type-Safety} &
  \textbf{Capabilities} &
  \textbf{Syntax} &
  \textbf{Ecosystem} \\
  \hline
  \endfirsthead
  \multicolumn{3}{@{}l}{} \\
  \hline
  \endhead
  \textbf{gettext}   & D & D & C & S \\
  \textbf{i18next}   & A & A & S & F \\
  \textbf{Paraglide} & S & B & B & F \\
  \textbf{Fluent}    & C & S & B & A \\
  \hline
  \textbf{Mulan} &
  S &
  D \textrightarrow{} S
  & B &
  D \textrightarrow{} A \\
  \hline
\end{longtable}
